\section{AHSME 1987}
        \begin{enumerate}
        	\item (\textit{AHSME 1987}) $(1+x^2)(1-x^3)$ equals


 $\text{(A)}\ 1 - x^5\qquad \text{(B)}\ 1 - x^6\qquad \text{(C)}\ 1+ x^2 -x^3\qquad \\  \text{(D)}\ 1+x^2-x^3-x^5\qquad \text{(E)}\ 1+x^2-x^3-x^6$


	\item (\textit{AHSME 1987}) A triangular corner with side lengths $DB=EB=1$ is cut from equilateral triangle ABC of side length $3$. 
The perimeter of the remaining quadrilateral is


 
                    \begin{figure}[H]
                        \centering
                        \includegraphics[width=0.5\textwidth]{Problem/AHSME_I_1987_2_Problem_Figure_1.png}
                    \end{figure}
                    
$\text{(A)} \ 6 \qquad  \text{(B)} \ 6\frac{1}{2} \qquad  \text{(C)} \ 7 \qquad  \text{(D)} \ 7\frac{1}{2} \qquad  \text{(E)} \ 8$


	\item (\textit{AHSME 1987}) How many primes less than $100$ have $7$ as the ones digit? (Assume the usual base ten representation) 


 $\text{(A)} \ 4 \qquad  \text{(B)} \ 5 \qquad  \text{(C)} \ 6 \qquad  \text{(D)} \ 7 \qquad  \text{(E)} \ 8$


	\item (\textit{AHSME 1987}) $\frac{2^1+2^0+2^{-1}}{2^{-2}+2^{-3}+2^{-4}}$ equals 


 $\text{(A)} \ 6 \qquad  \text{(B)} \ 8 \qquad  \text{(C)} \ \frac{31}{2} \qquad  \text{(D)} \ 24 \qquad  \text{(E)} \ 512$


	\item (\textit{AHSME 1987}) A student recorded the exact percentage frequency distribution for a set of measurements, as shown below. 
However, the student neglected to indicate $N$, the total number of measurements. What is the smallest possible value of $N$?


 
 \[\begin{tabular}{c c}\text{measured value}&\text{percent frequency}\\ \hline 0 & 12.5\\ 1 & 0\\ 2 & 50\\ 3 & 25\\ 4 & 12.5\\ \hline\ & 100\\ \end{tabular}\]
$\textbf{(A)}\ 5 \qquad \textbf{(B)}\ 8 \qquad \textbf{(C)}\ 16 \qquad \textbf{(D)}\ 25 \qquad \textbf{(E)}\ 50$


	\item (\textit{AHSME 1987}) In the $\triangle ABC$ shown, $D$ is some interior point, and $x, y, z, w$ are the measures of angles
in degrees. Solve for $x$ in terms of $y, z$ and $w$.


 
                    \begin{figure}[H]
                        \centering
                        \includegraphics[width=0.5\textwidth]{Problem/AHSME_I_1987_6_Problem_Figure_1.png}
                    \end{figure}
                    
$\textbf{(A)}\ w-y-z \qquad \textbf{(B)}\ w-2y-2z \qquad \textbf{(C)}\ 180-w-y-z \qquad \\ \textbf{(D)}\ 2w-y-z\qquad \textbf{(E)}\ 180-w+y+z$


	\item (\textit{AHSME 1987}) If $a-1=b+2=c-3=d+4$, which of the four quantities $a,b,c,d$ is the largest?


 $\textbf{(A)}\ a \qquad \textbf{(B)}\ b \qquad \textbf{(C)}\ c \qquad \textbf{(D)}\ d \qquad \textbf{(E)}\ \text{no one is always largest}$


	\item (\textit{AHSME 1987}) In the figure the sum of the distances $AD$ and $BD$ is


 
                    \begin{figure}[H]
                        \centering
                        \includegraphics[width=0.5\textwidth]{Problem/AHSME_I_1987_8_Problem_Figure_1.png}
                    \end{figure}
                    
$\textbf{(A)}\ \text{between 10 and 11} \qquad \textbf{(B)}\ 12 \qquad \textbf{(C)}\ \text{between 15 and 16}\qquad\\ \textbf{(D)}\ \text{between 16 and 17}\qquad \textbf{(E)}\ 17$


	\item (\textit{AHSME 1987}) The first four terms of an arithmetic sequence are $a, x, b, 2x$. The ratio of $a$ to $b$ is


 $\textbf{(A)}\ \frac{1}{4} \qquad \textbf{(B)}\ \frac{1}{3} \qquad \textbf{(C)}\ \frac{1}{2} \qquad \textbf{(D)}\ \frac{2}{3}\qquad \textbf{(E)}\ 2$


	\item (\textit{AHSME 1987}) How many ordered triples $(a, b, c)$ of non-zero real numbers have the property that each number is the product of the other two?


 $\textbf{(A)}\ 1 \qquad \textbf{(B)}\ 2 \qquad \textbf{(C)}\ 3 \qquad \textbf{(D)}\ 4 \qquad \textbf{(E)}\ 5$


	\item (\textit{AHSME 1987}) Let $c$ be a constant. The simultaneous equations

 \begin{align*}x-y = &\ 2 \\cx+y = &\ 3 \\\end{align*}

have a solution $(x, y)$ inside Quadrant I if and only if


 $\textbf{(A)}\ c=-1 \qquad \textbf{(B)}\ c>-1 \qquad \textbf{(C)}\ c<\frac{3}{2} \qquad \textbf{(D)}\ 0<c<\frac{3}{2}\\ \qquad \textbf{(E)}\ -1<c<\frac{3}{2}$


	\item (\textit{AHSME 1987}) In an office, at various times during the day the boss gives the secretary a letter to type, each time putting the letter on top of the pile in the secretary's in-box. When there is time, the secretary takes the top letter off the pile and types it. If there are five letters in all, and the boss delivers them in the order $1\ 2\ 3\ 4\ 5$, which of the following could not be the order in which the secretary types them?


 $\textbf{(A)}\ 1\ 2\ 3\ 4\ 5 \qquad \textbf{(B)}\ 2\ 4\ 3\ 5\ 1 \qquad \textbf{(C)}\ 3\ 2\ 4\ 1\ 5 \qquad \textbf{(D)}\ 4\ 5\ 2\ 3\ 1 \qquad \textbf{(E)}\ 5\ 4\ 3\ 2\ 1 \qquad$


	\item (\textit{AHSME 1987}) A long piece of paper $5$ cm wide is made into a roll for cash registers by wrapping it $600$ times around a cardboard tube of diameter $2$ cm, 
forming a roll $10$ cm in diameter. Approximate the length of the paper in meters. 
(Pretend the paper forms $600$ concentric circles with diameters evenly spaced from $2$ cm to $10$ cm.)


 $\textbf{(A)}\ 36\pi \qquad \textbf{(B)}\ 45\pi \qquad \textbf{(C)}\ 60\pi \qquad \textbf{(D)}\ 72\pi \qquad \textbf{(E)}\ 90\pi$


	\item (\textit{AHSME 1987}) $ABCD$ is a square and $M$ and $N$ are the midpoints of $BC$ and $CD$ respectively. Then $\sin \theta=$

                    \begin{figure}[H]
                        \centering
                        \includegraphics[width=0.5\textwidth]{Problem/AHSME_I_1987_14_Problem_Figure_1.png}
                    \end{figure}
                    
$\textbf{(A)}\ \frac{\sqrt{5}}{5} \qquad \textbf{(B)}\ \frac{3}{5} \qquad \textbf{(C)}\ \frac{\sqrt{10}}{5} \qquad \textbf{(D)}\ \frac{4}{5}\qquad \textbf{(E)}\ \text{none of these}$


	\item (\textit{AHSME 1987}) If $(x, y)$ is a solution to the system
$xy=6$ and $x^2y+xy^2+x+y=63$,
find $x^2+y^2$.


 $\textbf{(A)}\ 13 \qquad \textbf{(B)}\ \frac{1173}{32} \qquad \textbf{(C)}\ 55 \qquad \textbf{(D)}\ 69 \qquad \textbf{(E)}\ 81$


	\item (\textit{AHSME 1987}) A cryptographer devises the following method for encoding positive integers. First, the integer is expressed in base $5$. 
Second, a 1-to-1 correspondence is established between the digits that appear in the expressions in base $5$ and the elements of the set 
$\{V, W, X, Y, Z\}$. Using this correspondence, the cryptographer finds that three consecutive integers in increasing 
order are coded as $VYZ, VYX, VVW$, respectively. What is the base-$10$ expression for the integer coded as $XYZ$?


 $\textbf{(A)}\ 48 \qquad \textbf{(B)}\ 71 \qquad \textbf{(C)}\ 82 \qquad \textbf{(D)}\ 108 \qquad \textbf{(E)}\ 113$


	\item (\textit{AHSME 1987}) In a mathematics competition, the sum of the scores of Bill and Dick equalled the sum of the scores of Ann and Carol. 
If the scores of Bill and Carol had been interchanged, then the sum of the scores of Ann and Carol would have exceeded 
the sum of the scores of the other two. Also, Dick's score exceeded the sum of the scores of Bill and Carol. 
Determine the order in which the four contestants finished, from highest to lowest. Assume all scores were nonnegative.


 $\textbf{(A)}\ \text{Dick, Ann, Carol, Bill} \qquad \textbf{(B)}\ \text{Dick, Ann, Bill, Carol} \qquad \textbf{(C)}\ \text{Dick, Carol, Bill, Ann}\\ \qquad \textbf{(D)}\ \text{Ann, Dick, Carol, Bill}\qquad \textbf{(E)}\ \text{Ann, Dick, Bill, Carol}$


	\item (\textit{AHSME 1987}) It takes $A$ algebra books (all the same thickness) and $H$ geometry books (all the same thickness, 
which is greater than that of an algebra book) to completely fill a certain shelf. 
Also, $S$ of the algebra books and $M$ of the geometry books would fill the same shelf. 
Finally, $E$ of the algebra books alone would fill this shelf. Given that $A, H, S, M, E$ are distinct positive integers, 
it follows that $E$ is


 $\textbf{(A)}\ \frac{AM+SH}{M+H} \qquad \textbf{(B)}\ \frac{AM^2+SH^2}{M^2+H^2} \qquad \textbf{(C)}\ \frac{AH-SM}{M-H}\qquad \textbf{(D)}\ \frac{AM-SH}{M-H}\qquad \textbf{(E)}\ \frac{AM^2-SH^2}{M^2-H^2}$


	\item (\textit{AHSME 1987}) Which of the following is closest to $\sqrt{65}-\sqrt{63}$?


 $\textbf{(A)}\ .12 \qquad \textbf{(B)}\ .13 \qquad \textbf{(C)}\ .14 \qquad \textbf{(D)}\ .15 \qquad \textbf{(E)}\ .16$


	\item (\textit{AHSME 1987}) Evaluate
$\log_{10}(\tan 1^{\circ})+\log_{10}(\tan 2^{\circ})+\log_{10}(\tan 3^{\circ})+\cdots+\log_{10}(\tan 88^{\circ})+\log_{10}(\tan 89^{\circ}).$
$\textbf{(A)}\ 0 \qquad \textbf{(B)}\ \frac{1}{2}\log_{10}(\frac{\sqrt{3}}{2}) \qquad \textbf{(C)}\ \frac{1}{2}\log_{10}2\qquad \textbf{(D)}\ 1\qquad \textbf{(E)}\ \text{none of these}$


	\item (\textit{AHSME 1987}) There are two natural ways to inscribe a square in a given isosceles right triangle. 
If it is done as in Figure 1 below, then one finds that the area of the square is $441 \text{cm}^2$. 
What is the area (in $\text{cm}^2$) of the square inscribed in the same $\triangle ABC$ as shown in Figure 2 below?


 
                    \begin{figure}[H]
                        \centering
                        \includegraphics[width=0.5\textwidth]{Problem/AHSME_I_1987_21_Problem_Figure_1.png}
                    \end{figure}
                    
$\textbf{(A)}\ 378 \qquad \textbf{(B)}\ 392 \qquad \textbf{(C)}\ 400 \qquad \textbf{(D)}\ 441 \qquad \textbf{(E)}\ 484$


	\item (\textit{AHSME 1987}) A ball was floating in a lake when the lake froze. The ball was removed (without breaking the ice), 
leaving a hole $24$ cm across as the top and $8$ cm deep. What was the radius of the ball (in centimeters)?


 $\textbf{(A)}\ 8 \qquad \textbf{(B)}\ 12 \qquad \textbf{(C)}\ 13 \qquad \textbf{(D)}\ 8\sqrt{3} \qquad \textbf{(E)}\ 6\sqrt{6}$


	\item (\textit{AHSME 1987}) If $p$ is a prime and both roots of $x^2+px-444p=0$ are integers, then


 $\textbf{(A)}\ 1<p\le 11 \qquad \textbf{(B)}\ 11<p \le 21 \qquad \textbf{(C)}\ 21< p \le 31 \\ \qquad \textbf{(D)}\ 31< p\le 41\qquad \textbf{(E)}\ 41< p\le 51$


	\item (\textit{AHSME 1987}) How many polynomial functions $f$ of degree $\ge 1$ satisfy
$f(x^2)=[f(x)]^2=f(f(x))$  ?


 $\textbf{(A)}\ 0 \qquad \textbf{(B)}\ 1 \qquad \textbf{(C)}\ 2 \qquad \textbf{(D)}\ \text{finitely many but more than 2}\\ \qquad \textbf{(E)}\ \infty$


	\item (\textit{AHSME 1987}) $ABC$ is a triangle: $A=(0,0), B=(36,15)$ and both the coordinates of $C$ are integers. What is the minimum area $\triangle ABC$ can have?


 $\textbf{(A)}\ \frac{1}{2} \qquad \textbf{(B)}\ 1 \qquad \textbf{(C)}\ \frac{3}{2} \qquad \textbf{(D)}\ \frac{13}{2}\qquad \textbf{(E)}\ \text{there is no minimum}$


	\item (\textit{AHSME 1987}) The amount $2.5$ is split into two nonnegative real numbers uniformly at random, 
for instance, into $2.143$ and $.357$, or into $\sqrt{3}$ and $2.5-\sqrt{3}$. 
Then each number is rounded to its nearest integer, for instance, $2$ and $0$ in the first case above, 
$2$ and $1$ in the second. What is the probability that the two integers sum to $3$?


 $\textbf{(A)}\ \frac{1}{4} \qquad \textbf{(B)}\ \frac{2}{5} \qquad \textbf{(C)}\ \frac{1}{2} \qquad \textbf{(D)}\ \frac{3}{5}\qquad \textbf{(E)}\ \frac{3}{4}$


	\item (\textit{AHSME 1987}) A cube of cheese $C=\{(x, y, z)| 0 \le x, y, z \le 1\}$ is cut along the planes $x=y, y=z$ and $z=x$. How many pieces are there? 
(No cheese is moved until all three cuts are made.)


 $\textbf{(A)}\ 5 \qquad \textbf{(B)}\ 6 \qquad \textbf{(C)}\ 7 \qquad \textbf{(D)}\ 8 \qquad \textbf{(E)}\ 9$


	\item (\textit{AHSME 1987}) Let $a, b, c, d$ be real numbers. Suppose that all the roots of $z^4+az^3+bz^2+cz+d=0$ are complex numbers lying on a circle
in the complex plane centered at $0+0i$ and having radius $1$. The sum of the reciprocals of the roots is necessarily


 $\textbf{(A)}\ a \qquad \textbf{(B)}\ b \qquad \textbf{(C)}\ c \qquad \textbf{(D)}\ -a \qquad \textbf{(E)}\ -b$


	\item (\textit{AHSME 1987}) Consider the sequence of numbers defined recursively by $t_1=1$ and for $n>1$ by $t_n=1+t_{(n/2)}$ when $n$ is even 
and by $t_n=\frac{1}{t_{(n-1)}}$ when $n$ is odd. Given that $t_n=\frac{19}{87}$, the sum of the digits of $n$ is


 $\textbf{(A)}\ 15 \qquad \textbf{(B)}\ 17 \qquad \textbf{(C)}\ 19 \qquad \textbf{(D)}\ 21 \qquad \textbf{(E)}\ 23$


	\item (\textit{AHSME 1987}) In the figure, $\triangle ABC$ has $\angle A =45^{\circ}$ and $\angle B =30^{\circ}$. A line $DE$, with $D$ on $AB$ 
and $\angle ADE =60^{\circ}$, divides $\triangle ABC$ into two pieces of equal area. 
(Note: the figure may not be accurate; perhaps $E$ is on $CB$ instead of $AC$.) 
The ratio $\frac{AD}{AB}$ is


 
                    \begin{figure}[H]
                        \centering
                        \includegraphics[width=0.5\textwidth]{Problem/AHSME_I_1987_30_Problem_Figure_1.png}
                    \end{figure}
                    
$\textbf{(A)}\ \frac{1}{\sqrt{2}} \qquad \textbf{(B)}\ \frac{2}{2+\sqrt{2}} \qquad \textbf{(C)}\ \frac{1}{\sqrt{3}} \qquad \textbf{(D)}\ \frac{1}{\sqrt[3]{6}}\qquad \textbf{(E)}\ \frac{1}{\sqrt[4]{12}}$


\end{enumerate}